\chapter{Sim-to-Real Deployment}
\label{ch:sim_to_real}

\section{Overview}

The preceding chapters developed and evaluated an articulated tractor-trailer controller entirely within a 2D simulation environment. This chapter describes the transfer of that controller to a physical robotic platform, constituting the sim-to-real component of this thesis. Unlike a conventional model-based stack, the deployed system does not re-implement a hand-derived tracking law on the vehicle. Instead, the simulation-trained reinforcement-learning (RL) policies are executed directly on hardware through a thin bridging layer that reproduces the observation and action interfaces the policies were trained against. The objective is to demonstrate that the vehicle model, the observation definitions, and the learned control logic established in simulation can be ported to real hardware with minimal algorithmic redesign, and to characterize the engineering challenges encountered during that transfer.

The deployment is built on a deliberately stripped-down Autoware stack. The standard Autoware localization and perception pipelines are retained essentially unchanged, while the behavior, planning, and control layers are replaced by a custom RL bridge that runs the forward and reverse policies. Two pieces of supporting infrastructure are described in detail because they are specific to the articulated, sensor-light platform: an onboard LiDAR-based hitch-angle estimator that recovers the articulation angle without any trailer-mounted sensor, and a Smith predictor~\cite{smith1957closer} that compensates for actuator delay so that delay-free-trained policies behave correctly on a lagging physical actuator.

\section{Hardware Platform}

\begin{figure}[t]
  \centering
  \safegraphic[width=0.85\textwidth]{figures/hardware/platform.jpg}{Photo of the AgileX Hunter SE tractor coupled to the lab trailer through the revolute hitch.}
  \caption{The physical articulated platform: an AgileX Hunter SE tractor coupled to the lab trailer through a revolute hitch at the tractor rear axle.}
  \label{fig:hw_platform}
\end{figure}

\subsection{Tractor Vehicle}

The physical tractor is an AgileX Hunter SE, a compact electric wheeled unmanned ground vehicle with an Ackermann steering geometry. The relevant mechanical parameters are a wheelbase of $0.65$\,m, a wheel radius of $0.15$\,m, and a maximum steering angle of $0.436$\,rad ($25^\circ$). Vehicle actuation is commanded over a CAN bus on a 20\,ms cycle. Speed commands are encoded as 16-bit signed integers in units of 1\,mm/s (the commanded speed in m/s scaled by $1000$) and steering commands as 16-bit signed integers scaled by approximately $1320$ counts per radian ($\approx 0.8$\,mrad per count).

\subsection{Trailer}

The trailer is coupled to the tractor through a revolute hitch joint located marginally behind the tractor rear axle, which was added to the tractor-trailer model used for training. The lab trailer has a wheelbase of approximately $2.8$\,m and a body width of roughly $0.33$\,m. The articulation angle $\gamma$, defined as the difference between the tractor and trailer headings ($\gamma = \psi_\text{truck} - \psi_\text{trailer}$), is estimated onboard from the tractor's own LiDAR rather than measured by any trailer-mounted sensor. The procedure is detailed in Section~\ref{sec:hitch_estimation}. The estimate is published to the ROS2 network on the \texttt{/vehicle/trailer\_state} topic. The control bridge consumes the most recent published value, with no staleness gating currently applied. Two distinct angular limits appear in the stack and should not be conflated. The onboard estimator clamps the published articulation estimate to $|\gamma| \le 90^\circ$, whereas the deployment trailer model defines a tighter operational limit of $|\gamma| \le 80^\circ$ ($1.40$\,rad), itself below the $90^\circ$ jackknife condition used to terminate simulation episodes.

\subsection{Sensor Suite}

\begin{figure}[t]
  \centering
  \safegraphic[width=0.85\textwidth]{figures/hardware/sensor_suite.jpg}{Photo of the onboard sensor suite mounted on the tractor.}
  \caption{Onboard sensor suite: the front-mounted RoboSense Helios LiDAR, the Basler cameras, the u-blox GNSS antenna, and the WitMotion IMU.}
  \label{fig:hw_sensors}
\end{figure}

The vehicle is instrumented with the following sensors:

\begin{itemize}
    \item \textbf{LiDAR}: RoboSense Helios mechanical LiDAR mounted at the front of the tractor, publishing 3D point clouds on \texttt{/rslidar\_points}. This single sensor is reused for three purposes: NDT-based localization~\cite{biber2003ndt}, environment perception, and onboard hitch-angle estimation.
    \item \textbf{Cameras}: Six Basler Pylon industrial cameras (rear, rear-right, center, left, merging, and right) producing compressed color images. The cameras are not consumed by the LiDAR-based RL policies and are used for logging and monitoring only.
    \item \textbf{IMU}: WitMotion 6-DoF inertial measurement unit, with a gyroscope-bias estimator applied in the localization pipeline.
    \item \textbf{GNSS}: u-blox ZED-F9P receiver with RTK corrections, providing centimeter-level position estimates converted to MGRS coordinates for Autoware localization.
    \item \textbf{Wheel odometry}: Vehicle speed and steering-angle feedback recovered from CAN messages and fused with IMU data for dead-reckoning between GNSS and scan-matching updates.
\end{itemize}

\section{Software Architecture}

\subsection{Middleware and Framework}

The deployment stack runs on ROS2 Humble (Ubuntu 22.04)~\cite{macenski2022ros2} and is integrated into the Autoware autonomous-driving framework~\cite{kato2018autoware}. Autoware provides modular localization, perception, planning, and control pipelines whose interfaces are defined by standardized ROS2 message types. As introduced in the overview, this deployment keeps the standard sensing, localization, and perception modules and replaces only the planning and control layers with the learned RL bridge.

\subsection{Relationship to the Conventional Articulated Pipeline}

For context, Autoware's conventional approach to articulated vehicles is fully model-based. A truck-trailer mode manager routes goals to a freespace planner (a hybrid A\textsuperscript{*} search~\cite{dolgov2010hybridastar} that respects the non-holonomic constraints of the articulated vehicle), and the resulting reference trajectory is tracked by a model-predictive or LQR lateral controller. This model-based pipeline exists within the wider project and represents the standard, well-understood baseline for trailer control. The relevant control background is reviewed in Chapter~\ref{ch:literature_review}. The contribution of this thesis, however, is the learned controller, and the deployment described here therefore routes control through the RL bridge instead of the planner/MPC chain. The two are mutually exclusive at run time and are selected at launch. \Cref{fig:autoware_stack} shows which modules are retained and which are replaced.

\begin{figure}[t]
  \centering
  \resizebox{\textwidth}{!}{\begin{tikzpicture}[
  >=Stealth, font=\footnotesize, node distance=8mm and 12mm,
  block/.style={draw, rounded corners, align=center, inner sep=4pt, minimum height=10mm, text width=20mm},
  lbl/.style={midway, font=\scriptsize}
]
  \node[block] (sens) {Sensing\\{\scriptsize LiDAR, GNSS, IMU}};
  \node[block, right=of sens] (loc) {Localization\\{\scriptsize (NDT)}};
  \node[block, right=of loc] (perc) {Perception};
  \node[block, right=of perc, fill=black!8] (rl) {RL bridge\\{\scriptsize TD3 policy}};
  \node[block, right=of rl] (veh) {Vehicle\\{\scriptsize Hunter SE}};

  \draw[->] (sens) -- (loc);
  \draw[->] (loc) -- (perc);
  \draw[->] (perc) -- (rl);
  \draw[->] (rl) -- (veh);

  % Autoware retained group
  \begin{scope}[on background layer]
    \node[draw, dashed, rounded corners, fit=(loc)(perc), inner sep=5pt] (aw) {};
  \end{scope}
  \node[font=\scriptsize, above=1pt of aw] {Autoware (retained)};

  % replaced conventional pipeline
  \node[block, below=12mm of rl, fill=gray!12, text=gray!55, draw=gray!55] (conv)
        {freespace planner\\{\scriptsize + MPC / LQR}};
  \draw[->, gray!55, dashed] (perc.south) |- (conv.west);
  \draw[->, gray!55, dashed] (conv.east) -| (veh.south);
  \node[font=\scriptsize, gray!55, below=1pt of conv, align=center] {conventional pipeline (replaced)};
  \node[font=\scriptsize, above=1pt of rl] {learned};
\end{tikzpicture}
}
  \caption{Deployment software stack. Autoware's localization and perception are retained unchanged, while the planning and control layers (the conventional freespace planner and MPC/LQR tracker) are replaced by the learned RL bridge. The two control paths are mutually exclusive and selected at launch.}
  \label{fig:autoware_stack}
\end{figure}

\subsection{RL Bridge Node Architecture}

The learned control path comprises the following principal nodes:

\begin{description}
    \item[\texttt{lane\_reference\_node}] Loads the Lanelet2 map, selects the active lane from the ego odometry and the current goal, and publishes the lane centerline as a reference path together with a \texttt{drive\_enabled} flag and a \texttt{drive\_direction} flag (forward or reverse). This node supplies the path geometry against which the policy's tracking observation is computed.

    \item[\texttt{rl\_bridge\_node}] The core of the deployment. It loads the trained forward and reverse TD3 policies (Stable-Baselines3~\cite{raffin2021sb3} / PyTorch~\cite{paszke2019pytorch} checkpoints) and runs a 10\,Hz control loop. On each tick it assembles the policy observation in exactly the layout used during training, an 8-dimensional state vector (steering angle, articulation angle, tractor and trailer tracking errors, and look-ahead path curvature) concatenated with a 24-beam LiDAR range vector, queries the active policy deterministically, and converts its 2-dimensional action (a steering rate and a stop/go signal) into the steering-angle and speed commands expected by the vehicle interface, with the stop/go signal mapped onto the task's fixed speed or a halt exactly as the simulation interface layer does during training. Direction-dependent post-processing (action smoothing and kinematic rate scaling) is applied for the reverse policy, whose outputs are inherently more oscillatory than the forward policy's.

    \item[\texttt{hitch\_angle\_estimator}] Estimates the articulation angle from the tractor LiDAR and publishes it on \texttt{/vehicle/trailer\_state} (Section~\ref{sec:hitch_estimation}).

    \item[\texttt{initialpose\_shim}] Bridges the RViz \texttt{/initialpose} message to the Autoware pose-initializer service, allowing the operator to seed localization interactively.

    \item[\texttt{controller\_canbridge}] Translates \texttt{autoware\_\allowbreak control\_\allowbreak msgs/\allowbreak Control} commands into Hunter SE CAN frames and parses vehicle feedback (velocity, steering angle) back into ROS2 topics.
\end{description}

The forward and reverse policies are loaded from separate checkpoints and selected according to the \texttt{drive\_direction} flag, so that the same bridge handles both maneuvering regimes without retraining or reconfiguration.

\Cref{fig:rl_bridge} shows how these nodes connect.

\begin{figure}[t]
  \centering
  \resizebox{\textwidth}{!}{\begin{tikzpicture}[
  >=Stealth, font=\footnotesize, node distance=7mm and 18mm,
  block/.style={draw, rounded corners, align=center, inner sep=4pt, minimum height=9mm},
  lbl/.style={midway, font=\scriptsize, align=center}
]
  \node[block] (lane) {\texttt{lane\_reference\_node}};
  \node[block, below=of lane] (hitch) {\texttt{hitch\_angle\_estimator}};
  \node[block, below=of hitch] (loc) {Autoware localization\\(LiDAR NDT)};
  \node[block, right=22mm of hitch] (bridge) {\texttt{rl\_bridge\_node}\\10\,Hz TD3 policy};
  \node[block, right=of bridge] (can) {\texttt{controller\_canbridge}};
  \node[block, right=of can] (veh) {Hunter SE\\(CAN bus)};

  \draw[->] (lane)  -- (bridge) node[lbl,above]{reference path,\\ \texttt{drive\_direction}};
  \draw[->] (hitch) -- (bridge) node[lbl,above]{\texttt{/vehicle/}\\\texttt{trailer\_state}};
  \draw[->] (loc)   -- (bridge) node[lbl,below]{kinematic\\ state};
  \draw[->] (bridge) -- (can) node[lbl,above]{\texttt{Control}};
  \draw[->] (can)   -- (veh) node[lbl,above]{CAN\\ command};
  \draw[->] (veh.south) to[out=-90,in=-90,looseness=1.2]
        node[lbl,below]{velocity, steering feedback} (can.south);
\end{tikzpicture}
}
  \caption{RL bridge dataflow. The \texttt{lane\_reference\_node}, the LiDAR-based \texttt{hitch\_angle\_estimator}, and Autoware localization supply the policy observation to \texttt{rl\_bridge\_node}, which queries the TD3 policy at 10\,Hz and emits a control command that \texttt{controller\_canbridge} translates into Hunter SE CAN frames. Vehicle velocity and steering feedback return through the same bridge.}
  \label{fig:rl_bridge}
\end{figure}

\subsection{Hitch Angle Estimation}
\label{sec:hitch_estimation}

Because the trailer carries no articulation encoder and no trailer-mounted sensing, the articulation angle is recovered geometrically from the tractor's front LiDAR, which observes the front (and, at larger angles, the side) face of the trailer. The estimator operates on each incoming point cloud as follows.

\paragraph{Region of interest.} Points are first cropped to a box in the LiDAR frame that brackets the expected location of the trailer face. The lateral extent of this region is shifted to track the face as the trailer swings: using the current filtered estimate $\gamma$, the lateral offset is set to
\begin{equation}
    \Delta y = -\,D_\text{face}\,\sin\gamma,
    \label{eq:hitch_roi_shift}
\end{equation}
where $D_\text{face}$ is the approximate longitudinal distance from the hitch pivot to the face. This keeps the cropped points on the trailer face across the full range of articulation rather than only near $\gamma = 0$.

\paragraph{Rectangle fitting.} Rather than fitting a single planar normal, the estimator fits the orientation of the trailer body to the cropped points using the minimum-closeness rectangle criterion of Zhang et al.~\cite{zhang2017lshape}. For a candidate orientation $\theta$, the points are rotated into axis-aligned coordinates $(u, v)$ and scored by their closeness to the nearest edges of the bounding rectangle,
\begin{equation}
    C(\theta) = \sum_i \min\!\big(d^u_i,\, d^v_i\big)^2,
    \qquad
    d^u_i = \min\!\big(|u_i - u_\text{min}|,\, |u_i - u_\text{max}|\big),
    \label{eq:hitch_closeness}
\end{equation}
with $d^v_i$ defined analogously. The fitted orientation minimizes $C(\theta)$ augmented by a dimension-prior penalty that favors rectangles consistent with the known trailer width and length:
\begin{equation}
    \theta^\star = \arg\min_{\theta \in [0, \pi)} \;\big[\, C(\theta) + \lambda\, P_\text{dim}(\theta) \,\big].
    \label{eq:hitch_fit}
\end{equation}
This formulation is robust to whether only the front face, both the front and a side face (an ``L'' shape), or only a side face is visible, the dominant cause of failure for a naive single-face normal estimate at large articulation. The four $90^\circ$ heading candidates implied by $\theta^\star$ are disambiguated by choosing the one closest to the previous filtered estimate and most consistent with the dimension prior, yielding a trailer longitudinal direction $\hat{\mathbf{d}}_t = (\cos\theta^\star, \sin\theta^\star)$.

\begin{figure}[t]
  \centering
  \safegraphic[width=0.7\textwidth]{figures/hardware/hitch_lidar_fit.png}{LiDAR returns on the trailer face with the fitted rectangle used for articulation-angle estimation.}
  \caption{Onboard hitch-angle estimation from the tractor LiDAR: returns on the trailer face (points) with the fitted minimum-closeness rectangle, whose orientation gives the articulation angle $\gamma$.}
  \label{fig:hitch_fit}
\end{figure}

\paragraph{Articulation angle and filtering.} The raw articulation angle is the signed angle between the truck-forward axis $\hat{\mathbf{x}} = (1, 0)$ and $\hat{\mathbf{d}}_t$,
\begin{equation}
    \tilde\gamma = \operatorname{atan2}\!\big(\hat{x}_1 \hat{d}_{t,2} - \hat{x}_2 \hat{d}_{t,1},\; \hat{\mathbf{x}} \cdot \hat{\mathbf{d}}_t\big),
    \label{eq:hitch_raw_angle}
\end{equation}
to which a field-calibration sign $s$ and offset $\gamma_0$ are applied and the result clamped to the operational range, $\gamma = \operatorname{clip}(s\,\tilde\gamma + \gamma_0,\, -\gamma_\text{max},\, \gamma_\text{max})$. The sign convention is chosen so that the published value matches Autoware's $\gamma = \psi_\text{truck} - \psi_\text{trailer}$. The raw per-scan estimate is then passed through a constant-angular-velocity Kalman filter~\cite{kalman1960new} on the state $[\gamma,\ \dot\gamma]$, with Mahalanobis innovation gating to reject outlier scans, and a short moving-average is maintained in parallel as a diagnostic. The filtered angle and its rate populate the \texttt{hitch\_angle} and \texttt{hitch\_rate} fields of the published \texttt{TrailerState}.

This onboard, geometry-plus-filter estimator is the actual mechanism behind the \texttt{/vehicle/trailer\_state} signal. It is not an external measurement. Its noise characteristics, and the way they are handled, are revisited in the gap analysis (Section~\ref{sec:gap_hitch}).

\subsection{Trailer State Estimation}

Because the trailer carries no independent localization, its full pose is reconstructed from the tractor's localized kinematic state and the estimated articulation angle. Given the tractor's rear-axle position and heading $(\mathbf{p}_\text{truck},\, \psi_\text{truck})$ from \texttt{/localization/\allowbreak kinematic\_state}, the hitch-point position is
\begin{equation}
    \mathbf{p}_\text{hitch} = \mathbf{p}_\text{truck} + d_h \begin{bmatrix} \cos\psi_\text{truck} \\ \sin\psi_\text{truck} \end{bmatrix},
\end{equation}
where $d_h$ is the (small, near-zero) signed hitch offset relative to the tractor rear axle. The trailer heading follows directly from the articulation angle, $\psi_\text{trailer} = \psi_\text{truck} - \gamma$, and the trailer rear axle is located at
\begin{equation}
    \mathbf{p}_\text{trailer} = \mathbf{p}_\text{hitch} - L_t \begin{bmatrix} \cos\psi_\text{trailer} \\ \sin\psi_\text{trailer} \end{bmatrix},
\end{equation}
where $L_t \approx 2.8$\,m is the trailer wheelbase. This reconstruction mirrors the hitch-kinematics model derived in Chapter~\ref{ch:modeling}, confirming that the simulation formulation is directly reusable in the deployment stack, and it provides the trailer-referenced tracking errors that form part of the policy observation.

\subsection{Actuator-Delay Compensation: Smith Predictor}
\label{sec:smith_predictor}

The simulation trains the policies against an idealized actuator that responds instantaneously to commands, whereas the physical steering and drive systems exhibit a first-order lag of the order of $50$--$150$\,ms. A policy trained for instantaneous response will, when confronted with a delayed actuator, repeatedly over-correct and excite a lag-induced oscillation. To reconcile the two, a Smith-predictor-style compensator is implemented that lets the delay-free-trained policy act on a delay-compensated state estimate.

The compensator maintains two ``shadow'' instances of the same kinematic \texttt{StateSpace\allowbreak TractorTrailer} model used in simulation: a delay-free shadow with zero actuation time constant, and a delayed shadow whose time constant matches the measured hardware lag. On each control tick both shadows are advanced with the command currently propagating through the actuator pipeline. The discrepancy between the measured vehicle state and the delayed shadow forms a residual,
\begin{equation}
    \mathbf{r}_k = \mathbf{x}^\text{meas}_k - \mathbf{x}^\text{delayed}_k,
    \label{eq:smith_residual}
\end{equation}
which captures model mismatch beyond pure delay (integration scheme, ground-contact effects, friction). The state presented to the policy is the delay-free shadow corrected by this residual,
\begin{equation}
    \hat{\mathbf{x}}_k = \mathbf{x}^\text{free}_k + \mathbf{r}_k,
    \label{eq:smith_predict}
\end{equation}
so that the policy ``sees'' an effectively delay-free vehicle while the residual keeps the prediction anchored to reality. The shadows are re-synchronized to the measured state on reset and whenever a large state jump is detected. The compensator is implemented and enabled in the deployment launch configuration, though its shadow time constants still require per-platform tuning before the loop can be closed on hardware. Combined with the explicit steering-rate limit it forms the primary mitigation for the actuation gap discussed in Section~\ref{sec:gap_actuation}.

\Cref{fig:smith_predictor} shows the compensator structure.

\begin{figure}[t]
  \centering
  \begin{tikzpicture}[
  >=Stealth, font=\small, node distance=10mm and 18mm,
  block/.style={draw, rounded corners, align=center, inner sep=4pt,
                minimum height=9mm, text width=30mm},
  sum/.style={draw, circle, inner sep=1pt, minimum size=6mm, font=\footnotesize},
  lbl/.style={font=\scriptsize}
]
  \node[block] (policy) {Policy $\pi$\\(delay-free trained)};
  \node[block, right=of policy] (act) {Actuator + vehicle\\(physical, delayed)};
  \node[sum, right=of act] (s1) {$\pm$};
  \node[block, below=of act] (ds) {Delayed shadow\\$\tau=\text{lag}$};
  \node[block, below=of policy] (fs) {Delay-free shadow\\$\tau=0$};
  \node[sum, below=of s1] (s2) {$+$};

  \coordinate (u) at ($(policy.east)!0.5!(act.west)$);

  \draw[->] (policy) -- (act) node[lbl,midway,above]{$u$};
  \draw[->] (u) |- (fs);
  \draw[->] (u) |- (ds);
  \draw[->] (act) -- (s1) node[lbl,midway,above]{$x^{\mathrm{meas}}$};
  \draw[->] (ds) -| (s1) node[lbl,pos=0.75,right]{$x^{\mathrm{delayed}}$};
  \draw[->] (s1) -- (s2) node[lbl,midway,right]{$r$};
  \draw[->] (fs) -| (s2) node[lbl,pos=0.75,left]{$x^{\mathrm{free}}$};
  \draw[->] (s2) -| ($(policy.south)$) node[lbl,pos=0.25,below]{$\hat{x}=x^{\mathrm{free}}+r$};
\end{tikzpicture}

  \caption{Smith-predictor compensator. The command $u$ drives the physical vehicle and two kinematic shadows. The delayed shadow is subtracted from the measured state to form a residual $r$, which corrects the delay-free shadow to give the delay-compensated estimate $\hat{x} = x^{\mathrm{free}} + r$ presented to the policy.}
  \label{fig:smith_predictor}
\end{figure}

\paragraph{Learned dynamics model (in progress).} The Smith predictor's fidelity is bounded by the accuracy of the analytical kinematic shadow. Work is in progress to replace, or augment, that shadow with a neural-network approximation of the true system dynamics. A data pipeline converts recorded driving logs (rosbags) into time-aligned $10$\,Hz state--action--next-state datasets (deriving the trailer pose from the truck pose and the estimated hitch angle where it is not directly recorded) and a multilayer-perceptron dynamics model (\texttt{DynamicsMLP}) is trained on these datasets. The learned model is evaluated both by one-step prediction error against held-out logs and by closed-loop rollout divergence from the kinematic model, the latter serving to quantify and localize the sim-to-real dynamics gap. This component is not yet part of the live control loop and is reported here as ongoing work rather than a validated result.

\section{Sim-to-Real Gap Analysis}

Transferring a simulation-trained policy to physical hardware introduces discrepancies between the simulated and real-world dynamics. The following gap sources were identified and addressed during the deployment process.

\subsection{Vehicle Scale and Parameterization}
\label{sec:gap_scale}

The most fundamental difference between the simulation studies and the physical platform is one of scale. The simulation of \cref{ch:modeling} is parameterized for a full-size shunt truck (\cref{tab:vehicle_params}), whereas the deployment vehicle is the compact AgileX Hunter SE (wheelbase $0.65$\,m, maximum steering angle $0.436$\,rad), with under a quarter of the wheelbase and half the steering range of the reference terminal tractor. Rather than transferring the shunt-truck-scale policy onto the smaller vehicle and relying on it to generalize across scale, the deployment closes this gap by construction. The model structure of \cref{ch:modeling} is retained unchanged and re-parameterized to the Hunter's wheelbase, steering limit, and speed range, and a fresh policy is trained against that re-parameterized model. Because the governing equations are identical and only the parameter set changes, the observation and action interfaces the policy was trained against are preserved, and the scale mismatch is removed before deployment rather than being left as a residual gap to be compensated at run time.

\subsection{Actuation Dynamics}
\label{sec:gap_actuation}

The simulation model assumes instantaneous steering response, whereas the physical vehicle exhibits a first-order steering lag introduced by the CAN command latency and the Hunter SE's internal steering servo. This is the highest-impact discrepancy for a policy trained on instantaneous actuation. It is mitigated by the Smith-predictor compensator of Section~\ref{sec:smith_predictor}, which presents a delay-compensated state to the policy, and by a steering-rate limit applied to the commanded output, consistent with the $\dot{\delta}_\text{max}$ constraint used during simulation training.

\subsection{Hitch-Angle Sensing}
\label{sec:gap_hitch}

In simulation, the articulation angle $\gamma$ is available exactly from the vehicle state. On hardware it is estimated from LiDAR returns and is therefore subject to noise from sparse or partial returns on the trailer face, mis-segmentation of the region of interest, and momentary loss of the face at large articulation. These error modes are addressed within the estimator itself rather than by a downstream filter. The rectangle-fitting formulation (Equation~\ref{eq:hitch_fit}) is robust to partial face visibility, and the constant-angular-velocity Kalman filter with innovation gating rejects outlier scans and smooths the published signal. If the trailer face is lost entirely the estimator withholds new estimates, and the bridge continues from the last published value until the face is reacquired. No automatic disengagement on stale data is currently implemented.

\subsection{Path Distribution}

The policies are trained on smooth, spline-generated reference paths, whereas the deployment reference is the centerline of a Lanelet2 lane, which is piecewise-linear and may contain comparatively sharp vertices. Re-sampling the centerline does not fully reconcile the difference in path-curvature distribution between training and deployment, and this mismatch is a contributing factor to any residual tracking error observed on hardware. It is noted here as a characterization of the gap rather than something fully eliminated in the current deployment.

\subsection{Jackknife Prevention}

The simulation introduced a terminal episode condition when $|\gamma| > 90^\circ$ but provided no explicit recovery mechanism, since the episode simply resets. Rather than adding a separate runtime supervisor on the vehicle, jackknife resilience is instilled during training, through two complementary measures. First, a recovery-scenario reset distribution spawns a fraction of reverse episodes already in the near-jackknife region (articulation magnitudes of roughly $0.35$--$0.80$\,rad, i.e.\ $20$--$45^\circ$), so the policy is repeatedly exposed to high-articulation states its nominal training distribution rarely visits and must learn to steer the trailer back toward alignment. Second, an anti-jackknife override is applied within the training environment. Whenever the articulation exceeds approximately $0.7$\,rad ($40^\circ$), the policy's steering command is replaced by a maximum counter-steer that drives the articulation back toward zero. This prevents the policy from learning that oscillatory, near-jackknife behavior is survivable, since every dangerous excursion is corrected during training and chattering no longer pays off. No dedicated runtime recovery override is currently part of the deployed bridge. On hardware, jackknife avoidance therefore relies on the behavior learned through these measures together with the operational $|\gamma| \le 80^\circ$ limit.

\subsection{Localization Dependency}

The simulation provides ground-truth vehicle pose at every step. The hardware stack instead depends on the Autoware localization pipeline, which fuses RTK-GNSS, gyro-based odometry, and LiDAR NDT scan matching through an extended Kalman filter. GNSS outages or localization divergence therefore represent a failure mode that does not exist in simulation. The current deployment treats localization failure as an emergency-stop condition.

\section{High-Fidelity Geographic Simulation}
\label{sec:geo_sim}

The 2D Pygame environment of \cref{ch:modeling} is deliberately abstract. It captures the tractor-trailer kinematics and lane geometry but not the three-dimensional and photometric structure that the deployed perception and localization stack actually consumes. A complementary, higher-fidelity simulation environment was therefore built to exercise the full Autoware stack (LiDAR-based NDT localization, perception, and the RL bridge) inside a geographically faithful ``digital twin'' of a real operating site as a complement to the physical trials.

This environment is built on the Gazebo simulator~\cite{koenig2004gazebo}, extending the open-source Rudis Laboratories georeferenced world-generation toolchain~\cite{perseghetti2021rudisworlds}. Real locations are reconstructed as textured 3D meshes authored in Blender~\cite{blender}, whose physically based materials are baked into texture maps for efficient real-time rendering. Building on that toolchain, the target distribution-center site was rebuilt in Blender directly from public geographic data, with both layers imported through the Blender GIS plugin. Building footprints were pulled from OpenStreetMap~\cite{openstreetmap} and extruded to form the 3D building meshes, and these were draped over aerial orthoimagery from the ESRI World Imagery basemap, replacing the lower-resolution ground textures of the original toolchain with the higher-resolution ESRI imagery (\cref{fig:geo_sim_blender}). Because the reconstruction is driven entirely by the OpenStreetMap footprints and the ESRI imagery layer, the same procedure generates a Gazebo world for any real distribution center from its coordinates alone, rather than being tied to a single hand-authored site. Each world is anchored to GIS coordinates through the SDF \texttt{spherical\_coordinates} element, which fixes the simulation origin to a real latitude, longitude, and elevation. The templated world generator additionally computes the sun's azimuth and elevation from the site coordinates and a chosen UTC date and time, so that lighting and shadows match the real location and time of day. The \texttt{electrans} tractor and trailer models are then spawned into these worlds, so that the same articulated vehicle can be driven through a virtual replica of the target environment (\cref{fig:electrans_model,fig:geo_sim_gazebo}).

\begin{figure}[H]
    \centering
    \safegraphic[width=0.85\textwidth]{figures/ch6/6.5/blender_esri_reconstruction.png}{Blender reconstruction of the target distribution center: OpenStreetMap building footprints extruded over ESRI World Imagery aerial orthophotography, both imported through the Blender GIS plugin.}
    \caption{Reconstruction of the target distribution-center site in Blender. OpenStreetMap building footprints are extruded to form the 3D building meshes and draped over high-resolution aerial orthoimagery from the ESRI World Imagery basemap, both layers imported through the Blender GIS plugin. The same GIS-driven procedure reconstructs any real distribution center from its geographic coordinates.}
    \label{fig:geo_sim_blender}
\end{figure}

\begin{figure}[H]
    \centering
    \begin{subfigure}[b]{0.9\textwidth}
        \centering
        \safegraphic[width=\textwidth]{figures/gazebo/gazebo_dc_world.png}{Reconstructed distribution center exported to a georeferenced Gazebo world.}
        \caption{The reconstructed site exported to a georeferenced Gazebo world.}
        \label{fig:geo_sim_gazebo_world}
    \end{subfigure}
    \\[0.75em]
    \begin{subfigure}[b]{0.9\textwidth}
        \centering
        \safegraphic[width=\textwidth]{figures/ch6/6.5/gazebo_dc_loading_dock.png}{The electrans tractor-trailer spawned at a loading dock in the georeferenced Gazebo world.}
        \caption{The \texttt{electrans} tractor-trailer spawned at a loading dock within the world.}
        \label{fig:geo_sim_gazebo_vehicle}
    \end{subfigure}
    \caption{The georeferenced site running in Gazebo. The exported world (a) preserves the real building geometry and aerial ground texture, and the \texttt{electrans} articulated vehicle (b) is spawned into it so that the same policies evaluated in simulation can be exercised against the full Autoware sensing and localization stack.}
    \label{fig:geo_sim_gazebo}
\end{figure}

Because the world is georeferenced and rendered in three dimensions, this environment directly targets several of the sim-to-real gaps catalogued above. A LiDAR sensor model produces point clouds against real building and ground geometry, so NDT scan matching and the onboard hitch-angle estimator (\cref{sec:hitch_estimation}) can be exercised against realistic returns; the geographic anchoring lets the GNSS and map frames be reproduced consistently with the physical site; and the photorealistic rendering makes it possible to evaluate camera-based perception that the abstract 2D environment cannot represent. It is also the natural setting in which to prototype the infrastructure sensor nodes of \cref{sec:sensor_nodes}, since virtual LiDAR units can be placed at arbitrary surveyed locations in the georeferenced world. The environment and its GIS-driven world-generation pipeline are complete. What remains is the quantitative characterization of policy behavior within it, which is identified as future work in \cref{ch:conclusion}.

\subsection{Articulated Vehicle Model}
\label{sec:electrans_model}

The articulated vehicle spawned into the georeferenced worlds is the \texttt{electrans} tractor-trailer, shown in \cref{fig:electrans_model}. It is defined in the conventional ROS/Gazebo manner as a URDF assembled from \texttt{xacro} macros, rather than as a hand-written SDF model. The tractor shares its entire \texttt{xacro} description with the AgileX Hunter model used elsewhere in the project. The same parameterized macros define an Ackermann steering geometry (a pair of revolute front-steering links bounded by the maximum steering angle), continuous-rotation rear wheels coupled by a mimic joint, the link inertials, and the Gazebo friction and \texttt{ros2\_control} blocks. The two vehicles differ only in scale and in the body mesh. The Hunter's compact chassis mesh is replaced by a cab-over terminal-tractor body, and the model is scaled accordingly. Reusing a single \texttt{xacro} description in this way keeps the simulated tractor kinematically consistent with the physical Hunter platform on which the policies are deployed, while presenting the larger visual form of the terminal tractor this thesis targets.

The trailer is coupled to the tractor through a revolute hitch joint at the tractor rear axle, matching the hitch convention of the kinematic model in \cref{ch:modeling}. It is modeled as a single box body carried on one axle and reuses the same wheel meshes and inertial macros as the tractor, so that the complete articulated system is expressed through the shared macro library. Because the model is authored as standard xacro, it drops directly into both the Gazebo digital twin of this section and the wider Autoware simulation tooling without modification.

\begin{figure}[H]
    \centering
    \safegraphic[width=0.75\textwidth]{figures/ch6/6.5/electrans_model.png}{The electrans tractor-trailer model: a cab-over tractor coupled to a box trailer through a revolute hitch joint.}
    \caption{The \texttt{electrans} tractor-trailer model. The tractor reuses the AgileX Hunter \texttt{xacro} description (Ackermann front steering, continuous rear wheels, and \texttt{ros2\_control} blocks), rescaled and given a cab-over terminal-tractor body mesh. The box trailer is attached through a revolute hitch joint at the tractor rear axle.}
    \label{fig:electrans_model}
\end{figure}

\section{Infrastructure Sensor Nodes}
\label{sec:sensor_nodes}

A structural limitation of the onboard sensing arrangement is that the trailer occludes the region directly behind it. The tractor-mounted LiDAR cannot see past the trailer, and instrumenting the trailer itself is undesirable. It would add cabling, power, calibration, and a wireless link to an otherwise passive, interchangeable body, and would have to be repeated for every trailer. This is most acute during reverse maneuvering, precisely when visibility of the space behind the trailer matters most.

This limitation is addressed with off-vehicle infrastructure sensor nodes, fixed LiDAR units placed at surveyed vantage points around the operating area, each providing a bird's-eye view of obstacles (including the region behind the trailer) without any sensor being mounted on the trailer. Each node is calibrated into a shared global frame so that its detections can be expressed in the same coordinates as the onboard perception output and fused into Autoware's perception and occupancy-grid representation, yielding a persistent, occlusion-resilient view of the obstacles around the articulated vehicle that complements the limited onboard field of view.

This capability is implemented and running in the laboratory as the indoor-perception stack. Each infrastructure node runs Euclidean clustering on its own LiDAR returns to detect candidate objects. An offboard fusion pipeline transforms every node's detections into the shared global frame, merges detections that refer to the same physical object across nodes by their overlapping footprint polygons, and time-aligns the streams against each node's modeled latency before passing the merged detections to an extended-Kalman-filter multi-object tracker. The tracker emits a single stable, deduplicated obstacle list for the whole monitored area.

Because the laboratory sensor nodes and their perception stack run under ROS1 (Noetic) while the vehicle runs under ROS2 (Humble), the two are joined by a dedicated interface node. It subscribes to the ROS1 perception outputs and republishes the fused obstacle list over ROS2, where it is consumed by the vehicle's Autoware perception and occupancy-grid representation and thus reaches the RL bridge. Critically, this interface node runs offboard on the laboratory compute rather than on the vehicle. The vehicle's onboard Jetson therefore runs only ROS2, and never has to host a dual ROS1/ROS2 stack. The ROS1 dependency is confined entirely to the off-vehicle infrastructure, which both simplifies the embedded software image and keeps the version-sensitive ROS1 perception packages off the resource-constrained onboard computer.

In laboratory operation the infrastructure nodes recover the region directly behind the trailer that the onboard LiDAR cannot observe, extending obstacle coverage to [TODO]\% of the vehicle's immediate surroundings during reverse maneuvers, at an added end-to-end latency of [TODO]\,ms relative to onboard-only perception.

% TODO: Infrastructure sensor-node quantitative results (numbers only; system is implemented)
%   - Behind-trailer detection coverage vs. onboard-only baseline (percentage)
%   - End-to-end latency of the ROS1->ROS2 obstacle path (ms)
%   - Node placement / coverage figure for the deployed lab layout if available

\begin{figure}[H]
    \centering
    \resizebox{0.85\textwidth}{!}{\begin{tikzpicture}[
  >=Stealth, font=\footnotesize,
  node2/.style={draw, fill=black!70, text=white, rounded corners, inner sep=2pt, font=\scriptsize},
  box/.style={draw, rounded corners, align=center, inner sep=3pt, font=\scriptsize}
]
  % node positions
  \coordinate (A) at (-3.4,2.3);
  \coordinate (B) at (3.4,2.3);
  \coordinate (C) at (-3.4,-2.3);

  % coverage cones, each opening from a node toward the vehicle region
  \foreach \n in {A,B,C}{
    \fill[blue!8] (\n) -- (-2.4,0.55) -- (1.4,-0.55) -- cycle;
  }

  % onboard blind spot behind the trailer, covered by the infrastructure cones
  \fill[red!12] (-2.4,-0.55) rectangle (-1.5,0.55);
  \node[font=\scriptsize, red!60, align=center] at (-3.0,0) {onboard\\blind spot};

  % vehicle (top-down): trailer + tractor
  \draw[fill=black!12] (-1.5,-0.35) rectangle (0.6,0.35);
  \draw[fill=black!30] (0.6,-0.3) rectangle (1.35,0.3);
  \node[font=\scriptsize] at (-0.45,0) {trailer};
  \node[font=\scriptsize] at (0.98,-0.55) {tractor};

  % infrastructure LiDAR nodes
  \node[node2] at (A) {LiDAR};
  \node[node2] at (B) {LiDAR};
  \node[node2] at (C) {LiDAR};

  % offboard fusion and ROS2 delivery
  \node[box] (fus) at (3.6,-2.3) {offboard fusion\\ROS1 $\to$ ROS2};
  \draw[->,gray,dashed] (B) -- (fus);
  \draw[->,gray,dashed] (A) to[out=-90,in=150] (fus);
  \draw[->,gray,dashed] (C) to[out=-25,in=200] (fus);
  \draw[->] (fus) to[out=110,in=-20]
        node[font=\scriptsize, pos=0.55, right]{obstacle list (ROS2)} (1.35,-0.1);
\end{tikzpicture}
}
    \caption{Infrastructure sensor-node architecture: fixed off-vehicle LiDAR units give a bird's-eye view of obstacles around the articulated vehicle, covering the region behind the trailer that the onboard sensors cannot observe. Their detections are fused and tracked offboard and delivered to the vehicle over ROS2, so the onboard computer runs only ROS2.}
    \label{fig:sensor_nodes}
\end{figure}

\section{Deployment Results}

% TODO: Insert experimental results from hardware trials here.
% Suggested content:
%   - Qualitative description of observed behavior (forward tracking, reverse stability)
%   - Tracking error and hitch-angle statistics compared to simulation benchmarks
%   - Representative trajectory plots or photographs of the physical platform
%   - Any failure modes encountered that were not observed in simulation
% TODO: Describe the metric-capture methodology before the results are reported,
%       i.e. how tracking-error and hitch-angle statistics are logged from
%       rosbags, which topics are recorded, the time window/segment selection,
%       and how the reported summary statistics are computed.

Hardware validation trials were conducted on [TODO: describe test location and scenario]. The forward policy achieved [TODO: describe forward lane-following behavior], with a mean trailer cross-track error of [TODO]\,m and articulation maintained within $|\gamma| < [TODO]^\circ$ throughout nominal operation. Reverse maneuvering [TODO: describe outcome]. [TODO: report the effect of enabling the Smith predictor on oscillation and tracking once tuned.]

\begin{figure}[t]
  \centering
  \safegraphic[width=0.85\textwidth]{figures/hardware/deployment_trial.jpg}{Photo of the articulated platform during a hardware trial at the test site.}
  \caption{The articulated platform during a hardware trial. [TODO: note the test site and maneuver shown.]}
  \label{fig:deployment_trial}
\end{figure}

\section{Summary}

This chapter demonstrated that the tractor-trailer control architecture developed in simulation can be ported to a physical platform by executing the trained RL policies directly through a bridging layer, rather than by re-deriving an analytical controller on the vehicle. The deployment retains Autoware's standard localization and perception while replacing planning and control with an RL bridge that runs the forward and reverse policies at 10\,Hz over the same observation and action interfaces used in training. The principal engineering contributions are this RL bridge and its integration into a stripped-down Autoware stack, the onboard LiDAR hitch-angle estimator that recovers the articulation angle by rectangle fitting and Kalman filtering without any trailer-mounted sensor, and the Smith-predictor compensator (with an in-progress learned dynamics model) that reconciles delay-free-trained policies with a lagging physical actuator. The gap analysis confirms that the simulation captures the dominant dynamics relevant to controller design, while identifying actuation delay, hitch-angle sensing, path-distribution mismatch, and localization reliability as the primary points of divergence between the simulated and physical systems. Finally, the infrastructure sensor nodes, implemented and running in the laboratory, provide a route to occlusion-resilient perception behind the trailer without adding sensors to the trailer itself.
